\documentclass[10pt,a4paper]{amsart}
\usepackage[margin=19mm]{geometry}
\usepackage{amsmath,amssymb,mathtools}
\usepackage[scr=rsfs,cal=euler]{mathalfa}
\usepackage{xfrac}
\usepackage[unicode,hidelinks,bookmarks=false]{hyperref}
\newcommand{\ie}{i.e.\ }
\newcommand{\cf}{cf.\ }
\newcommand{\resp}{respectively\ }
\newcommand{\Subsection}{\subsection}
\providecommand{\nobreakdash}{\nobreak\hbox{-}}
\setlength{\emergencystretch}{2em}
\multlinegap0pt
\usepackage{bm}
\usepackage{bbm}
\usepackage{dsfont}
\usepackage{xspace}
\usepackage{cancel}
\usepackage{mathtools}
\usepackage{amsthm}
\makeatletter
\@ifundefined{thm}{\newtheorem{thm}{Thm}}{}
\@ifundefined{prop}{\newtheorem{prop}[thm]{Proposition}}{}
\makeatother
\def\Om{\Omega}
\def\bc{{\mathbb C}}
\def\bn{{\mathbb N}}
\def\br{{\mathbb R}}
\def\bz{{\mathbb Z}}
\def\cf{{\mathcal F}}
\def\ch{{\mathcal H}}
\def\eps{\varepsilon}
\def\gd{{\mathfrak D}}
\def\l{\lambda} \def\L{\Lambda}
\def\om{\omega} \def\Om{\Omega}
\def\t{\tau}
\title[Periodicity, type $II\_1$ factors and free Poisson laws in in]{Periodicity, type $II\_1$ factors and free Poisson laws in interacting Fock spaces}
\author{Vitonofrio Crismale, Yun Gang Lu, \'Eric Ricard}
\date{Source: arXiv:2606.18162v1 (CC BY 4.0)}
\begin{document}
\maketitle

\noindent\emph{Source and licence.} Periodicity, type $II_1$ factors and free Poisson laws in interacting Fock spaces. Version arXiv:2606.18162v1 (CC BY 4.0), distributed under the Creative Commons Attribution 4.0 International licence, as recorded in the arXiv licence metadata for this exact version and shown on the abstract page https://arxiv.org/abs/2606.18162v1. Full rights notice: licenses/arxiv-2606.18162-CC-BY-4.0.txt.\par\smallskip

\noindent\textit{Editorial scope.} Counted unit: the Prop whose source label is \texttt{2perfree} in the article (source0.tex lines 1294--1300, source label \texttt{2perfree}), delivered together with the article's own complete original proof of it (source0.tex lines 1301--1378, one \texttt{proof} environment of 78 lines). No mathematics is written, repaired or re-derived: statement and proof are the article's own text line for line. Required context retained uncounted: cr (lines 354--357); cr2 (lines 359--363); free1 (lines 371--377); 2ome01b (lines 412--424); moment1 (lines 435--453). Every definition, display and label of those lines is retained unchanged. Omitted from this package and not redistributed: 1--353, 358--358, 364--370, 378--411, 425--434, 454--1293, 1379--1455. Nothing inside a retained excerpt is omitted or altered: every retained body is delivered exactly as the article's lines are written, so no omission receipt of any kind accompanies this package. Citations: the retained excerpts carry no citation command at all, so no bibliography entry of the article is transcribed and no reference is redistributed. Reference adaptation: none was needed and none was made; every reference inside the delivered unit resolves inside it, so no unresolved reference or citation is produced. Printed slips kept literal: no printed word, symbol, hypothesis or number of the counted statement or of its proof was altered. Layout and class: the article's own class is \texttt{amsart} with packages including amsmath, amsthm, amsfonts, eucal, subcaption, tikz, pgfplots, epsfig, xy, graphicx, amssymb; the wrapper is the lane's pinned \texttt{amsart} editorial wrapper at 10pt on A4, which restates the theorem-like environments the retained text needs and adds the article's own macro definitions that the retained text needs (\texttt{\textbackslash Om}, \texttt{\textbackslash bc}, \texttt{\textbackslash bn}, \texttt{\textbackslash br}, \texttt{\textbackslash bz}, \texttt{\textbackslash cf}, \texttt{\textbackslash ch}, \texttt{\textbackslash eps}, \texttt{\textbackslash gd}, \texttt{\textbackslash l}, \texttt{\textbackslash om}, \texttt{\textbackslash t}), with extra wrapper packages (bm, bbm, dsfont, xspace, cancel, mathtools). The delivered numbering is the unit's own. Assets: no retained span contains a figure environment, a graphics-inclusion command, a PDF-page inclusion, a table or a TikZ picture; no image file is copied into this package, the package declares no asset dependency and no image file is redistributed with it. Licence: version arXiv:2606.18162v1 of this article is distributed under the Creative Commons Attribution 4.0 International licence, as recorded in the arXiv licence metadata for this exact version and shown on the abstract page https://arxiv.org/abs/2606.18162v1. A full rights notice is delivered at licenses/arxiv-2606.18162-CC-BY-4.0.txt. Characters: every retained excerpt is pure ASCII, so the delivered engine, XeLaTeX with its default Latin Modern text font, reports no missing character.
\begin{equation}
\label{cr}
a(f) a^{\dag} (g) = \langle f, g\rangle\,\om_{N+1},
\end{equation}

\begin{equation}
\label{cr2}
\om_{N+1} a^\dag(f) = a^\dag(f) \om_{N+2}, \qquad
\om_{N+1} a(f) = a(f) \om_N .
\end{equation}

\begin{prop}\label{free1}
For each $j\in\bn$, let $\gd^0_j$ be the linear span of
$$
\bigg\{(a_j^\dag)^{n_+}a_j^{n_-}\prod_{k=1}^{n_0}\om_{N+h_k} :n_0,n_-,n_+\in\bn,\, n_++n_-\in 2\bn,\, h_k\in\bz,\, k\in [n_0]\bigg\}.
$$
Then $\gd^0_j$ is a $*$-algebra containing the identity operator of $\cf(\ch,\l_n)$.
\end{prop}

\begin{align}
\label{2ome01b}
\t_\Om(s_j^n)=\langle\Omega,s_j^n\Omega\rangle
=
\begin{cases}
0 & \text{if $n$ is odd},\\
\displaystyle
\sum_{\eps\in\{-1,1\}^{2n}_+}
\langle\Omega,
a_j^{\eps(1)}\cdots a_j^{\eps(2m)}\Omega\rangle
& \text{if $n=2m$},
\end{cases}
\end{align}

\begin{prop}
\label{moment1}
For each $\eps\in\{-1,1\}^{2n}_+$ and each $j\in\bn$, let $\{l_h,r_h\}_{h=1}^n$ be the corresponding non crossing pair partition associated with $\eps$. If $f,g,f_1,\ldots, f_{2n}\in\ch$, then
\begin{equation}
\label{2ome02a}
\langle\Omega,
a^{\eps(1)}(f_1)\cdots a^{\eps(2n)}(f_{2n})\Omega\rangle
=
\prod_{h=1}^n \om_{2h-l_h}\langle f_{l_h},f_{r_h}\rangle,
\end{equation}
and
\begin{align}
\label{2ome02b}
\langle\Omega,
a(f) a^{\eps(1)}(f_1)\cdots a^{\eps(2n)}(f_{2n}) a^\dag(g) \Omega\rangle
=
\om_1\langle f,g\rangle \prod_{h=1}^n \om_{2h-l_h+1}\langle f_{l_h},f_{r_h}\rangle.
\end{align}
\end{prop}

\begin{prop}\label{2perfree}
For a 1-mode type interacting Fock space $\cf(\ch,\l_n)$, the following hold:
\begin{itemize}
\item[(i)] The family $(s_j)_j$ is $\t_\Om$- free independent in the vacuum state if and only if then the weights $(\omega_n)_{n\geq 1}$ are constant.
\item[(ii)]  The family $(s_j^2)_j$ is $\t_\Om$-free if and only if the weights $(\om_n)_n$ have period 2. Moreover, if $(\om_{n})$ has period 2 , the family $(\gd_j)_j$ is $\t_\Om$-free independent.
\end{itemize}
\end{prop}

\begin{proof}
  For (i), we note that, if the sequence $(\omega_n)_{n\geq 1}$ is constant, then
  the family interacting Fock space reduces to the full Fock space based on the Hilbert space $\lambda_1\ch$. In this case the $*$-algebras generated by annihilators are $\t_\Om$-free.\\
Conversely, we first recall that if $x$ and $y$ are free in the vacuum, then
  $\tau_\Om(x^*y^*yx)=\tau_\Om(x^*x)\tau_\Om(y^*y)$, and hence $\|yx\Omega\|=\|x\Omega\|\|y\Omega\|$. Since for each $n$
$$
s_2^n\Om=e_2^{\otimes n}
+\sum_{k=0}^{n-1}c_{n,k}e_2^{\otimes k}\,, \quad c_{n,k}\in\br
$$
the vectors $\Om,s_2\Om,\ldots,s_2^n\Om$ and $\Om,e_2,\ldots,e_2^{\otimes n}$
span the same subspace. Therefore, there exists a unique monic polynomial
$p_n$ of degree $n$ such that $p_n(s_2)\Om=e_2^{\otimes n}$. We choose $x=s_1$ and $y=p_n(s_2)$ %where $P\in \mathbb R[X]$ so that
  %P(s_2)\Om=e_{2}^{\otimes n}$,
  $n\geq 1$. We get $\|e_{2}^{\otimes n}e_1\|=\|e_{2}^{\otimes n}\|\|e_1\|$, that is $\lambda_{n+1}=\lambda_n\lambda_1$. Hence, $\omega_n=\lambda_n/\lambda_{n-1}=\lambda_1$ for all $n\geq 1$.\\
Let us prove (ii). We start with the second part. Suppose that the weights are 2-periodic. For $\t_\Om$-freeness of $(\gd_j)_j$, it is sufficient to prove that for any $n\in\bn$, $i_1\neq i_2\neq \cdots\neq i_n\in\bn$, one has
\begin{equation}\label{freen}
\t_\Om\big((x_1-\t_\Om(x_1)I)(x_2-\t_\Om(x_2)I)\cdots (x_n-\t_\Om(x_n)I)\big)=0\,,
\end{equation}
where for each $j$, $x_j=(a_{i_j}^\dag)^{n_{+,j}}a_{i_j}^{n_{-,j}}\prod_{k=1}^{n_{0,j}}\om_{N+h_k}$ is a linear generator
of the unital $*$-algebra $\gd^0_{i_j}$ appearing in Proposition \ref{free1}. \\
We first suppose that there exists at least a word $x_k$ in the chain
above such that $n_{+,k}+n_{-,k}=0$, and take $m$ the largest index
for which this holds. Thus, $x_m$ has the form
$\displaystyle{\prod_{l=1}^r\om_{N+h_l}}$, and for each $p>m$,
$\t_\Om(x_p)=0$ as $n_{+,p}+n_{-,p}>0$. Hence,
$\big(x_{m+1}-\t_\Om(x_{m+1})I\big)\cdots\big(x_{n}-\t_\Om(x_{n})I\big)=x_{m+1}\cdots
x_n$, and denote
$$
\displaystyle{M_+(m):=\sum_{j=m+1}^n n_{+,j}}\,, \quad \displaystyle{M_-(m):=\sum_{j=m+1}^n n_{-,j}}\,.
$$
If $M_+(m)<M_-(m)$, then $x_{m+1}\cdots x_n\Om=0$, as there are more annihilators than creators. Therefore, \eqref{freen} follows. Since in general $M_+(m)-M_-(m)\in 2\bz$, the case $M_+(m)\geq M_-(m)$ implies $M_+(m)-M_-(m)=2p$, for $p\in\bn$. As a consequence, \eqref{cr2} gives
\begin{align*}
&(x_{m}-\t_\Om(x_{m})I\big)\big(x_{m+1}-\t_\Om(x_{m+1})I\big)\cdots\big(x_{n}-\t_\Om(x_{n})I\big)\Om\\
&= \bigg(\prod_{l=1}^r\om_{N+h_l}-\prod_{l=1}^r\om_{h_l}\bigg)x_{m+1}\cdots x_n\Om\\
&=x_{m+1}\cdots x_n\bigg(\prod_{l=1}^r\om_{N+h_l+2p}-\prod_{l=1}^r\om_{h_l}\bigg)\Om\,.
\end{align*}
Here,
$$
\displaystyle{\bigg(\prod_{l=1}^r\om_{N+h_l+2p}-\prod_{l=1}^r\om_{h_l}\bigg)\Om=\bigg(\prod_{l=1}^r\om_{h_l+2p}-\prod_{l=1}^r\om_{h_l}\bigg)\Om}=0\,,
$$
as $\om_{l+2p}=\om_l$ for each $l\in\bn^*$. Thus, also in this case, \eqref{freen} follows.\\
Suppose now that for each $k=1,\ldots, n$, one has $n_{+,k}+n_{-,k}>0$. This gives $x_k-\t_\Om(x_k)I=x_k$, as $x_k$ has mean zero. Therefore, the chain $(x_1-\t_\Om(x_1)I)(x_2-\t_\Om(x_2)I)\cdots (x_n-\t_\Om(x_n)I)$ reduces to $x_1\cdots x_n$. Again \eqref{cr2} allows to shift each $\prod_{k=1}^{n_0}\om_{N+h_k}$ to the right of the above chain. Thus, if $F:\bz\rightarrow \bc$, one simply denotes
$$
x_1\cdots x_n=(a^\dag_{i_1})^{n_{+,1}}a_{i_i}^{n_{-,1}}\cdots (a^\dag_{i_i})^{n_{+,n}}a_{i_n}^{n_{-,n}}F(N)\,,
$$
where $F(N)$ is achieved through functional calculus.
Since $F(N)\Om=F(0)\Om$, \eqref{freen} is proved if one gets
$$
\t_\Om((a^\dag_{i_1})^{n_{+,1}}a_{i_i}^{n_{-,1}}\cdots (a^\dag_{i_i})^{n_{+,n}}a_{i_n}^{n_{-,n}})=0\,.
$$
If $n_{-,n}>0$, the thesis is got to. Therefore, we reduce ourselves to showing that
\begin{equation}\label{freen2}
\t_\Om((a^\dag_{i_1})^{n_{+,1}}a_{i_i}^{n_{-,1}}\cdots (a^\dag_{i_i})^{n_{+,n}})=0\,,
\end{equation}
which trivially follows if $n_{-,k}=0$ for each $k\leq n-1$. In contrast, take $m$ the largest index for which $n_{-,m}>0$. As in this case also $n_{+,m+1}>0$, \eqref{cr} implies \eqref{freen2}.\\
Now assume that $(s_j^2)_j$ is $\tau_\Om$-free independent. First, iterating the arguments in the proof of Proposition \ref{moment1}, for each $\eps\in\{-1,1\}^{2n}_+$ one has
\begin{equation}\label{dep2}
\t_\Om\big(a_ia_ia^{\eps(1)}_j \cdots a^{\eps(2n)}_j a_i^\dag a_i^\dag\big)=\om_1\om_2\prod_{h=1}^n \om_{2h-l_h+2}\,.
\end{equation}
Freeness yields $\t_\Om(s_i^2s_j^{2}s_i^2)=\t_\Om(s_i^4)\t_\Om(s_j^{2})$, which immediately entails $\om_1(\om_2\om_3+\om_1^2)=\om_1(\om_1\om_2+\om_1^2)$, and therefore $\om_3=\om_1$. An inspection of $\t_\Om(s_i^2s_j^{4}s_i^2)=\t_\Om(s_i^4)\t_\Om(s_j^{4})$ gives
$$
\om_1\om_2(\om_3\om_4+\om_3^2)+\om_1^2(\om_1\om_2+\om_1^2)=(\om_1\om_2+\om_1^2)^2\,,
$$
which implies $\om_2=\om_4$, as $\om_1=\om_3$. Suppose now that, for each $2\leq j\leq n+1$, $\om_j=\om_1$ if $j$ is odd, and $\om_j=\om_{2}$ if $j$ is even. By \eqref{2ome01b} and \eqref{dep2} it turns out
\begin{align*}
\t_\Om\big(s_i^2s_j^{2n}s_i^2\big)&=\t_\Om\big(a_ia_is_j^{2n}a_i^\dag a_i^\dag\big)+\t_\Om\big(a_ia_i^\dag s_j^{2n}a_ia_i^\dag\big)\\
&=\om_1\om_2\sum_{\eps\in\{-1,1\}^{2n}_+}\prod_{h=1}^n \om_{2h-l_h+2}\,\, +\om_1^2\sum_{\eps\in\{-1,1\}^{2n}_+}\prod_{h=1}^n \om_{2h-l_h}
\end{align*}
and
$$
\t_\Om(s_i^4)\t_\Om(s_j^{2n})=(\om_1\om_2+\om_1^2)\sum_{\eps\in\{-1,1\}^{2n}_+}\prod_{h=1}^n \om_{2h-l_h}
$$
Since by freeness $\t_\Om(s_i^2s_j^{2n}s_i^2)=\t_\Om(s_i^4)\t_\Om(s_j^{2n})$,
$$
\sum_{\eps\in\{-1,1\}^{2n}_+}\prod_{h=1}^n \om_{2h-l_h+2}=\sum_{\eps\in\{-1,1\}^{2n}_+}\prod_{h=1}^n \om_{2h-l_h}
$$
As above, one reduces to the partition with the maximum depth. Indeed, all other non crossing pair partitions involve weights already determined by the induction hypothesis, so that the only contribution depending on $\om_{n+2}$ arises in this case. Here, $\om_3\cdots \om_{n+1}\om_{n+2}=\om_1\cdots \om_{n-1}\om_n$. If $n$ is even, then the induction assumption gives $\om_{n+1}=\om_1$, and therefore $\om_{n+2}=\om_2$. The contrary happens if $n$ is odd.
\end{proof}

\end{document}
